% ======================================================================
% COMPLETE ENGLISH THESIS TEMPLATE AND SECTION-BY-SECTION WRITING GUIDE
% Single-file source. No external images, .bib file, or data are required.
% Compile: pdflatex Complete_English_Thesis_Template.tex (run THREE times).
% Overleaf: upload this file, select pdfLaTeX, then Recompile.
% Common packages from a standard TeX Live / MiKTeX installation are used.
%
% START HERE:
% 1. Read "How to Use This Template" in the compiled document.
% 2. Edit the metadata below.
% 3. Replace every \Draft{...} and all illustrative examples.
% 4. Adapt the chapter structure to your actual study and university rules.
% 5. Switch \showguidancefalse ONLY after you have read the guide.
% 6. Before submission, search the SOURCE for Draft, TODO, and template.
%
% IMPORTANT: This is a general research-thesis framework, not an approved
% university class. Institutional formatting and declaration rules prevail.
% Everything between \Guide{...} braces is instructional, not thesis prose.
% ======================================================================
\documentclass[12pt,a4paper,oneside]{report}

% --------------------- DOCUMENT SETTINGS -------------------------------
\usepackage[T1]{fontenc}
\usepackage[utf8]{inputenc}
\usepackage[english]{babel}
\usepackage{lmodern}
\usepackage[left=32mm,right=25mm,top=27mm,bottom=27mm]{geometry}
\usepackage{setspace}
\onehalfspacing
\usepackage{microtype}
\usepackage{amsmath,amssymb}
\usepackage{graphicx}
\usepackage{booktabs,tabularx,longtable,array}
\usepackage{enumitem}
\usepackage{xcolor}
\usepackage[most]{tcolorbox}
\usepackage{fancyhdr}
\usepackage{caption}
\usepackage{listings}
\usepackage[numbers,sort&compress]{natbib}
\usepackage{url}
\usepackage{hyperref}
\usepackage[nameinlink,noabbrev]{cleveref}
\usepackage{bookmark}

% --------------------- EDIT YOUR DETAILS -------------------------------
\newcommand{\ThesisTitle}{Your Full Thesis Title}
\newcommand{\ShortTitle}{Short Thesis Title}
\newcommand{\AuthorName}{Your Full Name}
\newcommand{\StudentID}{Your Student Number}
\newcommand{\DegreeName}{Master of Science in Your Discipline}
\newcommand{\DepartmentName}{Your Department}
\newcommand{\UniversityName}{Your University}
\newcommand{\SupervisorName}{Professor Firstname Lastname}
\newcommand{\CoSupervisorName}{Dr. Firstname Lastname}
\newcommand{\SubmissionDate}{Month Year}
\newcommand{\UniversityCity}{City, Country}

% --------------------- OPTIONAL COMPONENTS -----------------------------
% true = include; false = omit. Confirm local front-matter requirements.
\newif\ifshowguidance
\showguidancetrue       % Change to \showguidancefalse to hide guide boxes.
\newif\ifincludeapproval
\includeapprovalfalse   % Use the university's exact approval form if needed.
\newif\ifincludedeclaration
\includedeclarationtrue
\newif\ifincludededication
\includededicationfalse
\newif\ifincludepublications
\includepublicationsfalse
\newif\ifincludeglossary
\includeglossaryfalse

% --------------------- VISUAL STYLE ------------------------------------
\definecolor{ThesisBlue}{HTML}{173F5F}
\definecolor{GuideBlue}{HTML}{EDF4F9}
\definecolor{GuideBorder}{HTML}{9BB7CD}
\hypersetup{
  colorlinks=true,linkcolor=ThesisBlue,citecolor=ThesisBlue,urlcolor=ThesisBlue,
  pdftitle={\ThesisTitle},pdfauthor={\AuthorName},
  pdfsubject={Research thesis template and writing guide}
}
\setcounter{secnumdepth}{3}
\setcounter{tocdepth}{2}
\setlength{\headheight}{15pt}
\pagestyle{fancy}
\fancyhf{}
\fancyhead[L]{\small\nouppercase{\leftmark}}
\fancyfoot[C]{\thepage}
\renewcommand{\headrulewidth}{0.3pt}
\fancypagestyle{plain}{%
  \fancyhf{}\fancyfoot[C]{\thepage}\renewcommand{\headrulewidth}{0pt}}
\renewcommand{\chaptermark}[1]{\markboth{\thechapter.\ #1}{}}
\setlist[itemize]{topsep=4pt,itemsep=2pt,parsep=0pt}
\setlist[enumerate]{topsep=4pt,itemsep=2pt,parsep=0pt}
\captionsetup{font=small,labelfont=bf}
\renewcommand{\arraystretch}{1.18}
\newcolumntype{Y}{>{\raggedright\arraybackslash}X}
\lstset{basicstyle=\ttfamily\footnotesize,breaklines=true,
  frame=single,columns=fullflexible,keepspaces=true,showstringspaces=false}

% Guide boxes disappear in clean mode. Draft markers NEVER disappear:
% turning guidance off must not conceal missing thesis content.
\newcommand{\Guide}[1]{%
  \ifshowguidance
  \begin{tcolorbox}[enhanced,breakable,colback=GuideBlue,
    colframe=GuideBorder,boxrule=0.5pt,arc=1mm,
    title={Writing guide --- replace with your research},
    colbacktitle=ThesisBlue,coltitle=white,fonttitle=\small\bfseries,
    before skip=10pt,after skip=10pt]
    \small #1
  \end{tcolorbox}
  \fi}
\newcommand{\Draft}[1]{\par\noindent\textit{[TO WRITE: #1]}\par}
\newcommand{\FrontChapter}[1]{%
  \chapter*{#1}\phantomsection\addcontentsline{toc}{chapter}{#1}%
  \markboth{#1}{#1}}
\newcommand{\ChecklistItem}{\item[$\square$]}

\begin{document}
\hypersetup{pageanchor=false}

% ======================================================================
% TITLE PAGE
% ======================================================================
\begin{titlepage}
\centering
\vspace*{10mm}
{\Large\UniversityName\par}
\vspace{4mm}
{\large\DepartmentName\par}
\vspace{22mm}
{\Huge\bfseries\color{ThesisBlue}\ThesisTitle\par}
\vspace{15mm}
{\Large\AuthorName\par}
\vspace{3mm}
{\normalsize Student ID: \StudentID\par}
\vspace{14mm}
{\large A thesis submitted in partial fulfillment of the requirements\par
for the degree of\par}
\vspace{4mm}
{\large\bfseries\DegreeName\par}
\vspace{12mm}
{\normalsize Supervisor: \SupervisorName\par}
{\normalsize Co-supervisor: \CoSupervisorName\par}
% Delete the co-supervisor line if it does not apply.
\vfill
{\normalsize\UniversityCity\par\SubmissionDate\par}
\end{titlepage}

\pagenumbering{roman}
\hypersetup{pageanchor=true}

% ======================================================================
% TEMPLATE MANUAL (AUTOMATICALLY HIDDEN WHEN GUIDANCE IS OFF)
% ======================================================================
\ifshowguidance
\FrontChapter{How to Use This Template}
This document is a complete, editable framework for an English-language
research thesis, with particular support for engineering, computational,
experimental, and data-driven research. It is a writing scaffold: the blue
boxes describe what to write; italic bracketed markers identify content you
must supply. Neither is a research finding.

\section*{First decisions}
\begin{enumerate}
\item Obtain your institution's thesis handbook, approved title-page wording,
declaration, citation style, and submission requirements. These take priority
over this template. The margins, font, spacing, and chapter order here are
editable defaults, not universal standards.
\item Set the degree, title, author, university, supervisors, and date near
the top of the source. Delete the co-supervisor line if irrelevant.
\item Decide which chapters you need. A compact thesis may merge Chapters~2
and~3, incorporate Chapter~5 into methodology/results, or combine results
and discussion. A theoretical thesis may replace experiments with proofs.
A qualitative thesis may replace numerical metrics with an analytical and
coding framework. Do not add irrelevant experiments to fit the headings.
\item Use the research questions as the organizing thread: each question
needs an appropriate method, evidence, interpretation, and final answer.
\end{enumerate}

\section*{Build instructions}
Upload this single \texttt{.tex} file to Overleaf and select pdfLaTeX, or run
the following command locally three times to resolve the contents, lists,
citations, and cross-references:
\begin{quote}\ttfamily\small
pdflatex Complete\_English\_Thesis\_Template.tex
\end{quote}
No external figure or bibliography file is necessary for the initial build.
The reference at the end is an explicitly marked formatting placeholder.
Replace it with verified sources before writing a real literature review.

\section*{Guide mode and clean mode}
The source contains \texttt{\string\showguidancetrue}. Change it to
\texttt{\string\showguidancefalse} to hide this manual, all blue guide
boxes, and the writing-handbook appendix. The thesis chapter structure,
draft markers, illustrative table, equation, figure placeholder, and example
reference remain visible. You must replace or delete these manually. Clean
mode is a layout preview, not a submission-readiness check.

\section*{Suggested writing sequence}
Write methods and the record of your actual work first; then prepare results,
discussion, literature synthesis, and introduction. Write the conclusion and
abstract after your findings are stable. Revise the entire document so that
objectives, methods, results, and conclusions agree. There is no universal
required page count or reference count: use the scope agreed with your
supervisor and the evidence your argument needs.

\section*{Choosing a structure}
\noindent\begin{tabularx}{\textwidth}{@{}p{0.29\textwidth}Y@{}}
\toprule
Study type & Adaptation to this framework\\
\midrule
Computational / engineering & Retain theory, implementation, verification,
validation, numerical results, and uncertainty where relevant.\\
Experimental & Emphasize apparatus, calibration, sampling, measurement error,
controls, and reproducibility.\\
Quantitative empirical & Emphasize constructs, sampling, measurement,
statistical assumptions, effect sizes, and uncertainty.\\
Qualitative & Emphasize positionality, recruitment, data collection,
analytical approach, reflexivity, and supported interpretations.\\
Theoretical & Emphasize definitions, assumptions, propositions, proofs,
counterexamples, and implications.\\
\bottomrule
\end{tabularx}

\section*{Scope control}
This is a monograph-style framework. An article-based thesis may instead
require an integrative introduction, included papers, contribution statements,
and a synthesis chapter. Confirm the permitted format before reorganizing.
Keep only sections that help explain the research you actually completed.
\clearpage
\fi

% ======================================================================
% FRONT MATTER
% ======================================================================
\ifincludeapproval
\FrontChapter{Approval Page}
\Guide{Use the exact institutional approval/certification form. List the
committee, degree, defense date, and signatures only if required. Do not
insert simulated signatures or imply approval that has not occurred.}
\Draft{Insert the official approval page or approved wording.}
\fi

\ifincludedeclaration
\FrontChapter{Declaration of Originality and Contributions}
\Guide{Replace this page with the declaration required by your university.
State authorship, acknowledge collaborative contributions, identify any
previously published material, and disclose external assistance as required.
If translating an existing thesis, describe the relationship to the original
accurately; translation does not make it a new degree submission.}
\Draft{Insert the institution-approved declaration; do not treat this
instructional text as an official declaration.}
\vspace{12mm}
\noindent Name: \AuthorName\par
\vspace{6mm}
\noindent Signature: \rule{55mm}{0.4pt}\hfill Date: \rule{35mm}{0.4pt}
\fi

\ifincludededication
\FrontChapter{Dedication}
\Guide{Optional. Keep this brief and personal; omit it if not desired.}
\Draft{Write your dedication.}
\fi

\FrontChapter{Acknowledgements}
\Guide{Acknowledge supervisors, collaborators, technical support, funding
and grant identifiers, data providers, and personal support where appropriate.
Describe contributions accurately. Obtain consent before naming individuals
where needed. Keep scientific claims and results out of this section.}
\Draft{Write the acknowledgements.}

\FrontChapter{Abstract}
\Guide{Write a self-contained summary after completing the thesis. A useful
starting range is approximately 150--350 words, but follow your institution's
limit. Include: (1) context and problem; (2) the specific gap or question;
(3) approach, data, and scope; (4) the main actual results, quantified where
appropriate; and (5) the contribution and its limits. Avoid promises about
future work in place of findings. Usually omit citations, unexplained
abbreviations, equations, and information absent from the thesis.

\textbf{Sentence scaffold:} ``[Context] remains limited by [specific problem].
This thesis investigates [question]. Using [method] on [data/domain], the
study evaluates [comparison]. The results show [actual finding and
uncertainty]. These findings support [bounded implication].'' This is a
structural aid, not text to submit unchanged.}
\Draft{Write the final abstract, including actual findings.}
\vspace{4mm}
\noindent\textbf{Keywords:} \Draft{Provide roughly 4--6 searchable keywords,
or the number required by your institution.}

\ifincludepublications
\FrontChapter{Publications and Author Contributions}
\Guide{List related publications with accurate status (published, accepted,
submitted, or in preparation). Identify which thesis material overlaps and
your role in joint work. Check reuse permissions. Omit if irrelevant.}
\Draft{List relevant outputs and your contributions.}
\fi

\clearpage
\phantomsection
\addcontentsline{toc}{chapter}{Contents}
\tableofcontents
\clearpage
\phantomsection
\addcontentsline{toc}{chapter}{List of Figures}
\listoffigures
\clearpage
\phantomsection
\addcontentsline{toc}{chapter}{List of Tables}
\listoftables

\FrontChapter{List of Abbreviations}
\Guide{List only abbreviations actually used, normally alphabetically.
Define each at its first substantive occurrence as well. Remove the sample
entries if they are irrelevant; this list does not replace in-text definitions.}
\noindent\begin{tabularx}{\textwidth}{@{}lY@{}}
\toprule Abbreviation & Meaning\\\midrule
CFD & Computational fluid dynamics (illustrative entry)\\
ML & Machine learning (illustrative entry)\\
PDE & Partial differential equation (illustrative entry)\\
\bottomrule
\end{tabularx}
\Draft{Replace the sample abbreviations with those used in your thesis.}

\FrontChapter{Nomenclature}
\Guide{Include this list when mathematical symbols are used extensively.
Provide symbol, meaning, and unit; distinguish Latin/Greek symbols,
subscripts, and superscripts if the list is long. Use consistent notation.
The following entries are illustrative only.}
\noindent\begin{tabularx}{\textwidth}{@{}lYl@{}}
\toprule Symbol & Meaning & SI unit\\\midrule
$t$ & Time & s\\
$\mathbf{x}$ & Spatial position & m\\
$u$ & Generic dependent variable & Problem-dependent\\
$\theta$ & Model parameter vector & Problem-dependent\\
$\Omega$ & Spatial domain & Problem-dependent\\
\bottomrule
\end{tabularx}
\Draft{Replace the sample symbols, specify units, and add indices if needed.}

\ifincludeglossary
\FrontChapter{Glossary}
\Guide{Optional. Define specialized terms used repeatedly or differently
across disciplines. Avoid duplicating the abbreviation and symbol lists.}
\Draft{Add concise definitions and supporting sources where needed.}
\fi

\clearpage
\pagenumbering{arabic}

% ======================================================================
\chapter{Introduction}
\label{ch:introduction}
\Guide{Lead the reader from the research context to a precise problem, then
to the questions, approach, contribution, and scope. An introduction explains
why the study is needed; the literature review provides the detailed evidence.
End with a clear account of what this thesis will demonstrate.}

\section{Background and Context}
\Guide{Introduce the field, application, system, or phenomenon. Explain the
technical or scientific importance with verified evidence. Define essential
terms for a reader in your discipline who is unfamiliar with the exact topic.
Move from broad context to the particular setting; avoid an unrelated history.}
\Draft{Describe the context and why the problem matters.}

\section{Problem Statement}
\Guide{State what is currently unknown, inaccurate, inefficient, or poorly
understood. Identify the affected system or population and the consequences.
Separate the practical problem from the researchable knowledge gap. Make the
statement precise enough that the reader can judge whether you resolved it.}
\Draft{State the specific problem and its consequences.}

\section{Research Gap and Motivation}
\Guide{Summarize what existing approaches accomplish and what they leave
unresolved. Support the gap with sources analyzed in Chapter~2. Explain why
the limitation matters and how your study addresses it. Do not claim that
``no studies exist'' without a sufficiently documented search.}
\Draft{Identify the evidenced gap and motivation.}

\section{Research Aim and Objectives}
\Guide{Give one overarching aim, followed by specific, assessable objectives.
Use verbs such as develop, characterize, compare, estimate, validate, or
explain. Each objective should correspond to work you performed and evidence
you can present; avoid vague objectives such as ``understand everything.''}
\Draft{State the overall aim in one sentence.}
\begin{enumerate}[label=O\arabic*.]
\item \Draft{Define the first assessable objective.}
\item \Draft{Define the second assessable objective.}
\item \Draft{Define the third objective, or change the number to fit the study.}
\end{enumerate}

\section{Research Questions and, Where Appropriate, Hypotheses}
\Guide{Write questions answerable by your design. For hypothesis-driven work,
state the variables, expected relationship, and rationale; statistical null
and alternative hypotheses belong where they aid the analysis. Exploratory,
design, or qualitative research need not invent hypotheses.}
\begin{enumerate}[label=RQ\arabic*.]
\item \Draft{Write the first research question.}
\item \Draft{Write the second research question.}
\end{enumerate}

\section{Scope, Delimitations, and Assumptions}
\Guide{Specify the studied domain, population, geometry, operating range,
time period, datasets, and analysis boundaries as relevant. Delimitations
are deliberate scope choices; limitations are constraints on inference.
State assumptions and justify their intended range of validity.}
\Draft{Define what the study covers and the assumptions it relies on.}

\section{Contributions of the Thesis}
\Guide{List the actual contributions: new evidence, a method, a model,
dataset, implementation, theory, or evaluation. Distinguish a contribution
from a task completed. Describe novelty conservatively relative to prior
work, and identify what you personally contributed in collaborative work.}
\Draft{State the contributions and how the thesis supports each one.}

\section{Overview of the Approach}
\Guide{Give a short, high-level explanation of the study design and evidence.
Save implementation details and full equations for later chapters.}
\Draft{Summarize the research workflow in one or two paragraphs.}

\section{Organization of the Thesis}
\Guide{Describe the purpose of each retained chapter in one sentence.
Update this paragraph after changing the structure. Cross-references below
illustrate automatic numbering.}
\Draft{Introduce the literature review in \cref{ch:literature}, the framework
in \cref{ch:theory}, methods in \cref{ch:methods}, credibility checks in
\cref{ch:validation}, results in \cref{ch:results}, interpretation in
\cref{ch:discussion}, and conclusions in \cref{ch:conclusion}.}

% ======================================================================
\chapter{Literature Review}
\label{ch:literature}
\Guide{Build an argument for the study rather than listing paper summaries.
Organize by concepts, competing explanations, methods, or findings. Compare
what studies establish, their evidence, and their limitations. A narrative
review is not a systematic review unless systematic procedures were used.}

\section{Review Scope and Search Approach}
\Guide{State the themes, disciplines, period, databases, key search terms,
and inclusion/exclusion logic in proportion to the review's purpose. For a
systematic review, document the protocol, screening, eligibility, search
dates, appraisal, and synthesis procedures in sufficient detail. For a
narrative review, describe the selection approach honestly without implying
exhaustiveness. Include foundational and relevant recent work.}
\Draft{Explain how the literature was located and selected.}

\section{Key Concepts and Definitions}
\Guide{Define concepts and constructs consistently. Compare definitions
where disagreement affects your research. Identify the definition adopted
in this thesis and why it is appropriate.}
\Draft{Explain the essential concepts with appropriate sources.}

\section{Established Approaches and Baseline Methods}
\subsection{Foundational Approaches}
\Guide{Explain the central ideas, assumptions, evidence, and appropriate use
cases of established approaches. Discuss strengths as well as limitations.}
\Draft{Synthesize the foundational literature.}
\subsection{Recent Developments}
\Guide{Explain how later approaches address earlier limitations and what new
tradeoffs they introduce. Compare evaluation conditions before comparing
performance claims.}
\Draft{Synthesize relevant recent developments.}
\subsection{Alternative and Competing Approaches}
\Guide{Discuss credible alternatives, including simpler solutions. Explain
their relevance to your design and the rationale for your selected baselines.}
\Draft{Compare competing approaches and tradeoffs.}

\section{Critical Comparison of the Literature}
\Guide{Compare assumptions, data, evaluation settings, strengths, limitations,
and unresolved issues. A comparison table can help, but it needs an analytical
paragraph explaining the pattern. Do not copy authors' claims uncritically.}
\begin{table}[htbp]
\centering\small
\caption[Illustrative literature comparison]{Illustrative literature-review
table. Replace every cell with verified information and citations.}
\label{tab:literature}
\begin{tabularx}{\textwidth}{@{}p{0.15\textwidth}YYYY@{}}
\toprule Study & Approach & Evidence & Strength & Limitation\\\midrule
Study A & [Method] & [Data/test] & [Supported merit] & [Specific limit]\\
Study B & [Method] & [Data/test] & [Supported merit] & [Specific limit]\\
Your study & [Approach] & [Planned/actual evidence] & [Contribution] & [Scope]\\
\bottomrule
\end{tabularx}
\end{table}
\Draft{Explain the patterns and disagreements revealed by the comparison.}

\section{Synthesis and Refined Research Gap}
\Guide{Connect the literature directly to your questions. State what is
established, what remains uncertain, and what evidence your study adds.
Identify why the proposed methodology is appropriate for this gap.}
\Draft{Develop the literature-based rationale for the study.}

\section{Chapter Summary}
\Draft{Summarize the argument and transition to the theoretical framework.}

% ======================================================================
\chapter{Theoretical and Conceptual Framework}
\label{ch:theory}
\Guide{Explain the concepts, mechanisms, governing relations, or analytical
lenses needed to understand your methods. Merge this chapter with the
literature review or methodology if a separate chapter would repeat material.}

\section{System Description and Conceptual Model}
\Guide{Identify the system boundary, components, inputs, outputs, mechanisms,
and interactions. For empirical research, define constructs and proposed
relationships. For computational work, define the physical or mathematical
domain. Use a conceptual diagram only when it clarifies the relationships.}
\Draft{Describe the system or conceptual model.}

\section{Assumptions and Simplifications}
\Guide{List assumptions explicitly and explain their rationale and expected
effects. Distinguish assumptions imposed by theory from choices made for
tractability. Identify assumptions later tested by sensitivity analysis.}
\Draft{State and justify the main assumptions.}

\section{Governing Relations or Analytical Propositions}
\Guide{Present the equations, propositions, theoretical arguments, or
conceptual relationships on which the method depends. Define every symbol
and unit at first use, cite borrowed formulations, and state applicable
conditions. Put lengthy routine derivations in an appendix.}
\Draft{Present the governing relations or conceptual propositions.}

\subsection{Illustrative Equation Format}
\Guide{The following equation demonstrates formatting only. Replace it with
an equation relevant to your thesis, or delete this subsection. Define the
domain, operators, unknowns, and assumptions for any retained equation.}
\begin{equation}
  \mathcal{F}[u(\mathbf{x},t);\boldsymbol{\theta}]=0,
  \qquad (\mathbf{x},t)\in\Omega\times(0,T].
  \label{eq:governing}
\end{equation}
\noindent Here, $\mathcal{F}$ is a generic governing operator, $u$ an unknown
field, $\mathbf{x}$ position, $t$ time, $\boldsymbol{\theta}$ a parameter vector,
$\Omega$ a spatial domain, and $T$ a final time. These are illustrative
definitions and make no claim about your study.

\section{Boundary, Initial, and Admissibility Conditions}
\Guide{For mathematical/physical models, specify initial and boundary
conditions and their physical justification. For other designs, replace
this section with eligibility conditions or domain constraints. Remove it
when it does not apply.}
\Draft{Specify the conditions defining a well-posed or eligible problem.}

\section{Expected Mechanisms and Testable Implications}
\Guide{Explain how the framework motivates the questions, measurable
quantities, predictions, or interpretation. Avoid presenting expectations
as observations.}
\Draft{Connect the framework to the study design.}

\section{Chapter Summary}
\Draft{Summarize the framework used in the next chapter.}

% ======================================================================
\chapter{Research Methodology}
\label{ch:methods}
\Guide{Describe what you actually did and why it answers the research
questions. Provide enough detail for a knowledgeable reader to evaluate or
repeat the work. Distinguish preplanned procedures from exploratory changes.
Retain and rename the subsections appropriate to your study.}

\section{Research Design and Rationale}
\Guide{Identify whether the study is experimental, observational,
computational, theoretical, qualitative, mixed-methods, or design-oriented.
Explain why this design fits the question and which alternatives were
considered. Specify units of analysis and the conditions of comparison.}
\Draft{Describe and justify the research design.}

\section{Research Workflow}
\Guide{Describe the ordered stages from input collection to final analysis.
Identify decisions, dependencies, outputs, and quality checks. Include a
diagram when it adds clarity. The figure below requires no external file;
replace its framed placeholder with your own graphic.}
\begin{figure}[htbp]
\centering
\fbox{\parbox[c][35mm][c]{0.85\textwidth}{\centering
Replace with a research-workflow figure.\\[3mm]
Show inputs, processing or analysis, checks, and outputs.}}
% Example replacement:
% \includegraphics[width=0.9\textwidth]{figures/workflow.pdf}
\caption[Research workflow placeholder]{Research workflow placeholder.
Replace the figure and write a self-contained caption.}
\label{fig:workflow}
\end{figure}
\Draft{Explain the workflow and refer to \cref{fig:workflow} in the text.}

\section{Data, Materials, or Study Domain}
\subsection{Sources, Provenance, and Access}
\Guide{Describe datasets, specimens, documents, geometries, participants,
equipment, or simulation inputs. Give provenance, collection dates,
versions, permissions, licensing, and access conditions as applicable.
Distinguish observed, synthetic, simulated, and externally supplied data.}
\Draft{Identify the materials and their provenance.}
\subsection{Selection, Sampling, and Inclusion Criteria}
\Guide{Define the target population or parameter space, sampling method,
sample size rationale, exclusions, and actual retained sample. Explain
representativeness and potential selection bias. For simulations, define
case selection and parameter coverage.}
\Draft{Explain selection and the sample or case design.}
\subsection{Preprocessing and Quality Control}
\Guide{Document cleaning, normalization, missing-data handling, outlier
rules, filtering, transformations, and quality checks. Explain their order
and rationale. If learning from data, fit preprocessing only on training
data when required to prevent information leakage.}
\Draft{Describe preprocessing and quality-control decisions.}

\section{Variables, Constructs, and Operational Definitions}
\Guide{Define inputs, outcomes, predictors, controls, confounders, units,
scales, and measurement procedures as relevant. For qualitative research,
explain analytical concepts and how they guide attention without forcing
the data into unsupported categories.}
\Draft{Define the quantities and constructs used in the study.}

\section{Proposed Method, Model, or Analytical Procedure}
\subsection{Method Overview and Design Choices}
\Guide{Describe the method precisely and explain what is adopted, modified,
or original. Provide citations for existing methods. Justify important
choices by the question, constraints, and evidence.}
\Draft{Describe the core method and its rationale.}
\subsection{Detailed Formulation or Analysis Steps}
\Guide{Supply equations, algorithms, procedures, interview/coding protocols,
or proof strategy as applicable. Define objective functions, parameters,
constraints, update rules, and stopping criteria where relevant. Clarify
the relationship between components in coupled or multistage methods.}
\Draft{Provide sufficient detail to reconstruct the procedure.}
\subsection{Implementation or Execution}
\Guide{Record software and version numbers, libraries, hardware,
instrument settings, precision, seeds, solver tolerances, training settings,
and execution time as relevant. For qualitative work, record collection
procedures, researcher roles, coding process, and changes to the codebook.
Separate essential methodology from incidental code listings.}
\Draft{Describe the implementation or execution details.}

\section{Experimental, Simulation, or Evaluation Design}
\subsection{Cases, Conditions, and Comparisons}
\Guide{List the experimental conditions, simulation cases, or comparison
groups. Explain controls, randomization, blinding, repeated measures, or
parameter sweeps where relevant. State what is held constant.}
\Draft{Define the evaluation cases and conditions.}
\subsection{Baseline and Reference Methods}
\Guide{Select credible baselines, including a simple baseline if useful.
Describe their implementation and tuning. Use comparable inputs, data,
budgets, and evaluation rules, or explain unavoidable differences.}
\Draft{Specify baselines and the comparison protocol.}
\subsection{Data Partitioning and Independence}
\Guide{For predictive work, describe train/validation/test splits and
whether separation is by participant, time, location, or system. Prevent
related observations from leaking across partitions. If no data partition
is needed, explain independence or repeated-use considerations and rename
this subsection accordingly.}
\Draft{Describe data separation or independence assumptions.}
\subsection{Robustness, Sensitivity, and Ablation Plan}
\Guide{State how you vary uncertain inputs or important design choices.
For a multicomponent method, ablations isolate a component's contribution.
Use justified ranges and repeated runs where needed; do not add ablations
when there are no separable components.}
\Draft{Define the additional analyses and their purpose.}

\section{Evaluation Criteria and Analysis Plan}
\subsection{Metrics and Success Criteria}
\Guide{Define every metric, units, aggregation, direction of improvement,
and limitations. Explain why it answers the question. State pre-specified
success criteria if there were any; label criteria chosen after inspecting
results as retrospective.}
\Draft{Define the evaluation criteria and justify their use.}
\subsection{Statistical or Interpretive Analysis}
\Guide{For quantitative analysis, describe estimands, models/tests,
assumptions, uncertainty intervals, effect sizes, repeated observations,
and multiplicity handling as appropriate. Distinguish exploratory analyses.
For qualitative work, describe interpretation, coding, reflexivity, evidence
selection, and disagreement resolution; use reliability statistics only
when consistent with the chosen methodology.}
\Draft{Specify and justify the analysis procedure.}

\section{Ethics, Permissions, and Data Management}
\Guide{Report approvals, consent, confidentiality, security, licensing,
conflicts of interest, and data retention when relevant. Include approval
identifiers only if obtained. If formal review is not applicable, give a
brief accurate explanation under your institution's rules. Describe access
restrictions instead of promising to share data you cannot release.}
\Draft{Describe applicable ethics and data-management arrangements.}

\section{Reproducibility and Research Records}
\Guide{Identify the code version, environment, configurations, data versions,
seeds, run records, and instructions needed to repeat the study. State which
materials are available, where, and under what restrictions. Include stable
repository or archive identifiers only when real.}
\Draft{Explain how the work can be inspected and reproduced.}

\section{Methodological Limitations}
\Guide{Identify design constraints, possible biases, measurement weaknesses,
and assumptions that affect inference. Explain mitigations and residual
effects. Save the implications of observed limitations for the discussion.}
\Draft{State the main limitations of the design.}

\section{Chapter Summary}
\Draft{Summarize the design and its connection to the research questions.}

% ======================================================================
\chapter{Verification, Validation, and Credibility Checks}
\label{ch:validation}
\Guide{Retain a separate chapter when these checks are a substantial part
of the research; otherwise integrate them into methods/results. For numerical
models, verification examines whether the implementation solves the stated
model correctly; validation assesses agreement with suitable evidence for
the intended use. Neither establishes universal correctness. In empirical
or qualitative work, use the measurement-validity or trustworthiness checks
appropriate to the design instead of irrelevant numerical tests.}

\section{Purpose and Acceptance Criteria}
\Guide{State which aspect of credibility each check evaluates and why the
reference is suitable. Identify acceptance thresholds, their rationale,
and whether they were specified before observing the results.}
\Draft{Define the credibility questions and checks.}

\section{Verification or Internal Quality Checks}
\subsection{Analytical, Synthetic, or Controlled Tests}
\Guide{For code/models, compare against analytical solutions, manufactured
solutions, conservation laws, controlled inputs, or independently checked
calculations where applicable. For other studies, describe instrument,
coding, transcription, or data-integrity checks. State what each check
can and cannot establish.}
\Draft{Report the relevant internal checks and their outcomes.}
\subsection{Resolution, Convergence, or Stability}
\Guide{When applicable, test spatial/time resolution, sampling resolution,
optimizer convergence, tolerance sensitivity, or stability across repeated
runs. State the tested range, criterion, observed change, and chosen setting.
Low training loss alone does not demonstrate predictive accuracy.}
\Draft{Report applicable convergence or stability evidence.}

\section{Validation or External Credibility}
\Guide{Compare with independent measurements, trusted references,
out-of-sample observations, established benchmarks, or other appropriate
evidence. Do not use the same evidence both to tune a model and to claim
independent validation without explaining the dependence. For qualitative
work, describe contextually suitable credibility practices.}
\Draft{Present the external comparison or credibility evidence.}

\section{Uncertainty and Error Sources}
\Guide{Distinguish measurement, sampling, numerical, parameter, and model
uncertainty as relevant. Explain estimation and propagation methods and
the uncertainty not captured by your analysis. Separate random variation
from systematic bias where evidence allows.}
\Draft{Describe the uncertainty and its effect on confidence in the results.}

\section{Chapter Summary and Suitability for the Main Study}
\Draft{State what the checks support, any failures, and the justified
operating range of the subsequent analysis.}

% ======================================================================
\chapter{Results}
\label{ch:results}
\Guide{Present the evidence in the order needed to answer the questions.
Clearly separate main findings, secondary analyses, and exploratory findings.
Describe observations here; develop the broader explanation in Chapter~7,
unless a combined results-and-discussion structure suits your discipline.
Report negative and inconclusive results when relevant.}

\section{Overview of the Completed Study}
\Guide{Report the actual analyzed sample, cases, successful/failed runs,
exclusions, and deviations from the planned procedure. Summarize relevant
data or case characteristics without repeating the full methodology.}
\Draft{Describe what was actually analyzed.}

\section{Findings for Research Question 1}
\subsection{Primary Evidence}
\Guide{State the question, present the relevant figure/table/theme/proof,
and describe the main finding. Give units, denominators, number of cases
or runs, and uncertainty as applicable. Evidence must support the precision
and scope of the claim.}
\Draft{Present the main evidence addressing RQ1.}
\subsection{Supporting Comparisons}
\Guide{Present the comparisons that help assess the primary finding.
Report magnitude and uncertainty, not only significance labels. For
qualitative analysis, link interpretations to appropriately contextualized
extracts and anonymized identifiers.}
\Draft{Present the supporting evidence for RQ1.}

\section{Findings for Research Question 2}
\Guide{Use a parallel structure where it improves comparison. Add, remove,
or rename question-specific sections to match the actual research questions.}
\Draft{Present the evidence addressing RQ2.}

\section{Comparison with Baseline or Reference Approaches}
\Guide{Use the protocol defined in methods. Report comparable measures,
uncertainty, and relevant costs, such as runtime, data, or resources. If a
comparison is not directly fair, state the difference beside the result.}
\Draft{Present baseline/reference comparisons.}

\section{Sensitivity, Robustness, and Component Analyses}
\Guide{Present effects of changing assumptions, settings, inputs, or model
components. Distinguish stable findings from condition-dependent ones.
Show variability across repeated runs when relevant.}
\Draft{Report the planned additional analyses and any exploratory extensions.}

\section{Failure Cases and Unexpected Findings}
\Guide{Document cases where the method fails, findings conflict with
expectations, or evidence is inconclusive. Identify the conditions and
frequency when possible. Do not quietly remove difficult cases.}
\Draft{Report relevant failure cases and unexpected findings.}

\section{Summary of Key Findings}
\Guide{Summarize only findings established in this chapter. Organize them
by question or objective so that the discussion can interpret them directly.}
\Draft{Summarize the central findings without introducing new evidence.}

% ======================================================================
\chapter{Discussion}
\label{ch:discussion}
\Guide{Explain what the findings mean, why they may have occurred, how they
relate to prior work, and how far the conclusions extend. Make the chain
from evidence to interpretation explicit. Avoid repeating results verbatim.}

\section{Answers to the Research Questions}
\Guide{For each question, state the supported answer and cite the relevant
result sections. Distinguish supported, partially supported, unresolved,
and contradicted expectations. Explain uncertainty and alternative readings.}
\Draft{Interpret the findings as answers to each research question.}

\section{Interpretation of Mechanisms and Patterns}
\Guide{Connect observed patterns to theory or plausible mechanisms. Separate
demonstrated mechanisms from hypotheses. Consider competing explanations,
confounding, and measurement effects; correlation alone does not establish
causation.}
\Draft{Explain the findings and assess alternative explanations.}

\section{Comparison with Previous Research}
\Guide{Identify agreement and disagreement with specific studies. Explain
differences in data, context, assumptions, or methods that may account for
them. Use fair comparisons and avoid treating every difference as superiority.}
\Draft{Position the findings within the literature.}

\section{Scientific, Methodological, and Practical Implications}
\subsection{Scientific or Theoretical Implications}
\Guide{Explain which understanding, model, or theoretical expectation is
refined and what evidence supports that refinement.}
\Draft{State the scientific or theoretical implications.}
\subsection{Methodological Implications}
\Guide{Explain what other researchers can learn about study design,
measurement, modeling, analysis, or evaluation from this work.}
\Draft{State the methodological implications.}
\subsection{Practical Implications and Tradeoffs}
\Guide{Explain who could use the findings, under which conditions, and at
what cost or risk. State operational requirements and tradeoffs relevant
to real use. Do not claim deployment readiness based only on limited tests.}
\Draft{Explain the practical implications within the evidence-supported scope.}

\section{Limitations and Threats to Validity}
\Guide{Discuss limitations revealed by the completed study: selection,
measurement, implementation, analysis, confounding, model assumptions,
researcher interpretation, and limited coverage as applicable. Explain the
direction or likely consequence of each limitation rather than merely
listing generic weaknesses. Avoid equating transparency with elimination
of bias.}
\Draft{Assess the limitations and their consequences for interpretation.}

\section{Generalizability and Transferability}
\Guide{Define the populations, geometries, operating regimes, contexts,
or settings to which the findings may reasonably apply. State where
additional evidence is needed. For qualitative research, discuss the
contextual detail that supports a reader's assessment of transferability.}
\Draft{Define the supported reach of the findings.}

\section{Chapter Summary}
\Draft{Summarize the interpretation and transition to the final conclusions.}

% ======================================================================
\chapter{Conclusions and Future Work}
\label{ch:conclusion}
\Guide{Close the argument clearly. Answer the stated problem and questions,
identify the contribution, and acknowledge the most consequential limits.
Do not introduce new experiments, datasets, or results in this chapter.}

\section{Research Summary}
\Draft{Briefly restate the problem, aim, scope, and approach.}

\section{Principal Conclusions}
\Guide{Write concise conclusions supported by results. Map them to the
objectives and research questions. Where useful, include the actual result
magnitude and conditions; avoid copying the abstract or making wider claims.}
\Draft{State the main conclusions in question/objective order.}

\section{Contributions and Significance}
\Draft{Explain the final contributions and why they matter, within scope.}

\section{Recommendations}
\Guide{Include evidence-based recommendations for researchers, practitioners,
or decision-makers only when justified. State the conditions needed for
their use. Omit this section if the thesis does not support recommendations.}
\Draft{Provide justified recommendations, or remove this section.}

\section{Future Research}
\Guide{Prioritize a small set of specific next studies. Link each to a
limitation or open question and explain the method or evidence needed.
Avoid generic suggestions such as ``use more data'' without a rationale.}
\Draft{Describe feasible, prioritized directions for future work.}

\section{Closing Statement}
\Draft{End with a brief, evidence-supported statement of the thesis's value.}

% ======================================================================
% REFERENCES: self-contained starter bibliography, not real research.
% Replace the placeholder entry and cite verified sources using \cite{key}.
% For a large thesis, consider BibTeX or biblatex with an external .bib file.
% If adopting an external manager, REMOVE the thebibliography environment.
% Follow the university's required reference style instead of assuming
% this numeric example is mandatory. Do not mix bibliography systems.
% ======================================================================
\clearpage
\renewcommand{\bibname}{References}
\phantomsection
\addcontentsline{toc}{chapter}{References}
\begin{thebibliography}{99}
\bibitem{replace-this-reference}
\textbf{FORMATTING PLACEHOLDER --- NOT A REAL SOURCE.}
Replace with verified author(s), full title, journal/publisher, year,
volume/issue/pages as relevant, and DOI or stable URL where available.
\end{thebibliography}

% ======================================================================
\appendix
\chapter{Supplementary Derivations and Technical Details}
\label{app:derivations}
\Guide{Include supporting material that would interrupt the main argument:
lengthy derivations, additional proofs, parameter details, supplementary
checks, or expanded technical documentation. Cite the appendix where it
becomes relevant in the main text. Necessary reasoning must remain in the
main thesis; appendices should not hide essential evidence.}
\section{Additional Derivations or Proofs}
\Draft{Provide supplementary derivations or remove this section.}
\section{Additional Model or Instrument Specifications}
\Draft{Add the technical specifications needed for inspection or replication.}

\chapter{Supplementary Results and Study Materials}
\label{app:supplement}
\Guide{Include additional tables, figures, questionnaires, interview guides,
coding frameworks, protocols, or case descriptions as relevant. Respect
permissions and confidentiality. Give each item a useful title and context;
reference it from the main text.}
\section{Supplementary Results}
\Draft{Add supporting results that do not need to appear in the main chapter.}
\section{Study Instruments or Coding Materials}
\Draft{Add relevant instruments or protocols, or remove this section.}

\chapter{Reproducibility and Availability Record}
\label{app:reproducibility}
\Guide{Use this appendix as a practical replication record. Include only
items relevant to the study. Do not publish passwords, private keys,
restricted personal data, or confidential material in configurations.}
\section{Software, Hardware, and Version Record}
\Draft{Record exact software/library versions, operating system, hardware,
repository version/commit, and instrument identifiers as applicable.}
\section{Data and Configuration Record}
\Draft{Record data identifiers/versions, directory structure, configurations,
parameter values, seeds, preprocessing outputs, and access restrictions.}
\section{Execution or Analysis Instructions}
\Draft{Provide ordered commands or procedural steps to reproduce the main
outputs, with expected inputs and outputs and approximate resource needs.}
\section{Availability Statement}
\Draft{State where materials are accessible, their license or conditions,
and any justified restrictions. Do not insert non-existent repository links.}

% ======================================================================
% WRITING HANDBOOK: guidance mode only; not part of a submitted thesis.
% ======================================================================
\ifshowguidance
\chapter{Writing, Formatting, and Submission Handbook}
\label{app:handbook}
This appendix belongs to the template, not the finished thesis. It disappears
when \texttt{\string\showguidancefalse} is selected. The other appendices
remain, so remove irrelevant ones manually.

\section{Chapter Planning and Traceability}
Before drafting at length, build a question-to-evidence map. Replace the
following planning prompts in your own notes; the map need not appear in
the thesis unless it helps the reader.

\noindent\begin{tabularx}{\textwidth}{@{}p{0.13\textwidth}YYY@{}}
\toprule Question & Method & Evidence & Final answer\\\midrule
RQ1 & What procedure addresses it? & Which result supports it? & What can
you conclude, and within what scope?\\
RQ2 & What comparison or analysis? & Which table, figure, theme, or proof?
& What remains uncertain?\\
\bottomrule
\end{tabularx}

Choose chapter length by argumentative need. The literature review should
justify the gap, methods should enable evaluation, results should establish
the evidence, and discussion should explain its meaning. Avoid adding pages
or references simply to reach an arbitrary number.

\section{Writing Clear Academic English}
\begin{itemize}
\item Give each paragraph one purpose: state its point, support it with
evidence or reasoning, and connect it to the argument.
\item Use precise terms consistently. Prefer concrete statements such as
``the error decreased under these conditions'' to unsupported adjectives
such as ``excellent'' or ``revolutionary.''
\item Use present tense for established relationships and what a figure
shows; past tense is often suitable for completed procedures and observed
results. Follow your discipline's conventions consistently.
\item Distinguish observation, inference, and speculation. Match verbs to
the evidence: ``suggests'' or ``is consistent with'' may be more appropriate
than ``proves.''
\item Use active voice when it clarifies who did what. Use first person only
as accepted in your discipline and institution.
\item Define abbreviations on first use. Keep spelling conventions,
capitalization, decimal notation, symbols, and terminology consistent.
\item Explain why an equation, figure, or table matters; do not make the
reader infer its role from its presence alone.
\end{itemize}

\section{Results Versus Discussion}
\noindent\begin{tabularx}{\textwidth}{@{}YY@{}}
\toprule Results & Discussion\\\midrule
What was observed, estimated, established, or found? & What does the
finding mean, and why might it occur?\\
What is the magnitude, variability, context, or supporting evidence? & How
does it relate to theory, prior work, and alternative explanations?\\
Which planned and exploratory analyses were completed? & What conclusions
and implications are justified, and what limits them?\\
\bottomrule
\end{tabularx}

If your field favors a combined chapter, alternate evidence and interpretation
in clearly organized subsections. Preserve the distinction in the prose.

\section{Figures and Tables}
\begin{itemize}
\item Refer to every figure/table in the text, explain its purpose, and place
it near its first substantive discussion when practical.
\item Use informative captions. Define abbreviations, sample counts,
uncertainty bars, colors, and conditions needed to interpret the display.
State source/adaptation and permissions when applicable.
\item Use legible labels and units. Use vector graphics for line plots and
diagrams where practical; inspect raster figures at final print size.
\item Make visual comparisons fair: comparable axes, clearly stated scales,
and transparent normalization. Do not let truncated axes mislead.
\item Tables should report comparable precision and units. Avoid presenting
the same information repeatedly as a table, figure, and long paragraph.
\item In this template, table captions are above tables and figure captions
are below figures. Change this if your institution specifies otherwise.
\end{itemize}

\section{Equations, Units, and Notation}
Define symbols before or immediately after first use. Number equations
that you refer to later; use automatic references, as in
\cref{eq:governing}. Distinguish scalars, vectors, matrices, operators,
and indices consistently. State units and check dimensional consistency
where appropriate. Punctuate displayed equations as parts of sentences.
Use a single notation for a quantity unless a change is explicitly explained.

\section{Citations, References, and Evidence}
\begin{itemize}
\item Cite claims, methods, equations, datasets, software, and borrowed
ideas where attribution is needed. Prefer reading the original source.
\item Check that a citation supports the exact claim and scope stated.
Do not infer a paper's conclusion from its title or an automated summary.
\item Keep a source log: claim, source, relevant location, and notes on
limitations. Distinguish your paraphrase from copied wording in notes.
\item Paraphrase by expressing the idea accurately in your own structure;
changing a few words is insufficient. Mark direct quotations and provide
locations in the style required by your institution.
\item Verify author names, titles, publication status, year, venue, and
identifiers. Never invent references. Cite software/data versions when
relevant to reproducibility.
\item Include all cited sources and remove uncited bibliography entries
unless your institution permits a separate wider bibliography. There is
no universal minimum number of references for a thesis.
\end{itemize}
The starter uses \texttt{natbib} with a manually maintained numeric
bibliography. After adding a real entry with a unique key, cite it using
\texttt{\string\cite\{your-key\}}. For author--year or an external reference
manager, select the appropriate institutional style and adapt the citation
setup consistently. Do not retain the example placeholder reference.

\section{Translation, Editing, and AI Assistance}
If writing from a thesis originally drafted in another language, preserve
the meaning of technical statements, assumptions, units, figure labels,
and findings. Create a terminology list before translating. Check the
English against the original chapter by chapter, especially equations and
numerical values. Keep publication and submission history accurate.

Follow your institution's rules on language editing, translation, and AI
assistance. Describe assistance truthfully when a disclosure is required,
and check all generated wording, calculations, and citations yourself.
Do not allow a translation or editing tool to invent results, sources, or
stronger conclusions. If disclosure belongs in a specific institutional
form, use that form rather than improvising a declaration here.

\section{Practical LaTeX Editing}
\begin{itemize}
\item Put a descriptive \texttt{\string\label\{...\}} after each chapter,
section, equation, and caption you will reference. Use
\texttt{\string\cref\{...\}} rather than typing a number manually.
\item Escape special characters in ordinary prose: use
\texttt{\string\%}, \texttt{\string\&}, \texttt{\string\_},
\texttt{\string\#}, and \texttt{\string\$}.
\item Add external images with \texttt{\string\includegraphics} only after
uploading the corresponding files. Use relative paths for portability.
\item Write mathematics in math mode. Avoid inserting screenshots of
equations when editable equations are feasible.
\item For a larger project, split chapters into files using
\texttt{\string\input\{chapters/intro\}}. Keep one main document
and a consistent folder structure.
\item Recompile after changes to headings, labels, or citations. Resolve
undefined references, missing files, and overflowing lines before delivery.
\end{itemize}

\section{Revision Workflow}
\begin{enumerate}
\item \textbf{Argument pass:} Check that the problem, gap, questions,
methods, evidence, and conclusions form one coherent argument.
\item \textbf{Evidence pass:} Verify every numerical value, claim, citation,
table entry, unit, and figure against its source or analysis record.
\item \textbf{Structure pass:} Remove repetition, irrelevant headings,
unsupported assertions, and background that does not serve the study.
\item \textbf{Language pass:} Improve clarity, paragraph flow, terminology,
tense, grammar, and consistency without changing scientific meaning.
\item \textbf{Presentation pass:} Review the compiled PDF page by page,
including title page, contents, lists, equations, page breaks, and appendices.
\item \textbf{Supervisor pass:} Obtain feedback on scope, contribution,
interpretation, and compliance; keep a record of revisions.
\end{enumerate}

\section{Final Submission Checklist}
\begin{itemize}[leftmargin=1.6em]
\ChecklistItem The university's required template, title wording, margins,
font, spacing, page numbering, declarations, and file rules are satisfied.
\ChecklistItem All author, supervisor, institution, degree, date, and title
metadata are final; optional pages have been approved or removed.
\ChecklistItem Every research question has an evidence-supported answer,
and objectives, abstract, results, and conclusions agree.
\ChecklistItem All template instructions, sample entries, illustrative
objects, draft markers, and irrelevant chapters/appendices are replaced
or removed. Turning guide mode off is insufficient by itself.
\ChecklistItem Figures/tables are readable, numbered, cited in text, and
properly attributed. Equations, units, symbols, and notation are consistent.
\ChecklistItem Citations resolve, references are verified, and permissions,
authorship, ethics, and assistance disclosures are accurate where applicable.
\ChecklistItem Data, code, instrument, configuration, and availability
statements match what actually exists and can be accessed.
\ChecklistItem The compiled PDF has no undefined references, clipped
material, unintended blank pages, or misleading visual comparisons.
\ChecklistItem The final abstract describes completed findings; the
conclusion introduces no new results and states the important limits.
\ChecklistItem The source, required assets, final PDF, and research records
are backed up and labeled with a clear version/date.
\end{itemize}
\fi

\end{document}
